\subsection{Baire Category Theorem}
\lecdate{(check date) Lec 12 - Feb 24 (Week 7)}

motivation: what does it mean for a pt in a space to be ``generic"?
for instance, in $\bR$, most pts are in some sense, probably irrational.
note:
\begin{equation*}
	\bar{\bQ}=\bR \qquad \bar{\bQ^{c}}=\bR \qquad \bar{\bQ\cap\bQ^{c}}=\eset
\end{equation*}
so density alone isnt a good test.

\begin{lm}
	let $A,B\subseteq M$ be open dense. then $A\cap B$ is open dense.
\end{lm} \

\begin{pf}[source=Primary Source Material]
	fix $x\in M$ and $\ep>0$.
	since $A$ dense, $\ex a\in A$ s.t. $a\in B(x,\ep)$.
	since $B(x,\ep)\cap A$ open, $\ex r>0$ s.t. $B(a,r)\subseteq B(x,\ep)\cap A$.
	since $B$ dense, $\ex b\in B(a,r)$.
	then $b\in A\cap B\cap B(x,\ep)$ as needed.
\end{pf}
we can now state bct.

\begin{thm}[title=Baire Category Theorem]
	sps $M$ complete. let $G_{n}$ be a seq of open dense subsets of $M$. then:
	\begin{equation*}
		G \ = \ \bigcap_{n=1}^{\infty}G_{n}
	\end{equation*}
	is dense.
\end{thm}
sets of this form have many names, such as:
\begin{itemize}
	\item thick.
	\item dense $G_{\delta}$ sets, i.e. a ctbl intersection of open sets
\end{itemize}

\begin{crll}
	if $M$ complete and:
	\begin{equation*}
		M=\bigcup_{n=1}^{\infty}F_{n}
	\end{equation*}
	where $F_{n}$ closed, then at least one $F_{n}$ has non-empty interior.
\end{crll} \

\begin{pf}[source=Primary Source Material]
	if $M$ is as above, then $\eset=\bigcap_{n}F_{n}^{c}$.
	exercise: $F_{n}^{c}$ is dense iff $F_{n}$ has no interior.
\end{pf}
we now prove bct.

\begin{pf}[source=Primary Source Material]
	let $x\in M$, $\ep>0$. wts $\ex p\in G\cap B(x,\ep)$.
	we'll construct a cauchy seq inductively, w pts $p_{n}$ and radii $r_{n}$ s.t.
	$B(p_{1},2r_{1})\subseteq B(x,\ep)$, and:
	\begin{equation*}
		B(p_{n},2r_{n})\subseteq B(p_{n-1},r_{n-1})\cap \bigcap_{k=1}^{n-1}G_{k}
		\ =: \ B
	\end{equation*}
	to start, pick $p_{1}=x$ and $r_{1}=\ep/2$.
	now, for any $n$, $\bigcap_{k=1}^{n-1}G_{k}$ open dense. thus, $\ex y \in B$.
	since $B$ is open, $\ex r>0$ s.t. $B(y,r)\subseteq B$.
	thus, take $p_{n}=y$ and $r_{n}=\min(r/2,1/n)$.
	since $r_{n}\sto0$, $p_{n}$ is cauchy, thus $p_{n}\sto p$.
	since $p_{k}\in B(p_{n},r_{n})$ for all $k\geq n$, we have:
	\begin{equation*}
		p\in\bar{B(p_{n},r_{n})}\subseteq B(p_{n},2r_{n})
	\end{equation*}
	since $p\in B(p_{n},2r_{n})$ for all $n$, then clearly $p\in B(x,\ep)$ and $p\in G_{n}$ $\fa n$
	as needed.
\end{pf}
kinda missed what he said but sth abt ill-behaved fn's.

\begin{prop}
	$\ex$ a thick $A\subseteq C_{0}([a,b])$ s.t. $\fa f\in A$:
	\begin{enumerate}
		\item the left/right derivs dne for any $x\in[a,b]$
		\item $f$ is not monotone on any $I=(c,d)$.
	\end{enumerate}
\end{prop} \

\begin{pf}[source=Primary Source Material]
	define:
	\begin{gather*}
		\Delta_{h}f(x) \ := \ \frac{f(x+h)-f(x)}{h} \\
		R_{n} \ := \ \set{f:\fa x\in[a,b-1/n],\ex h<1/n,\abs{\Delta_{h}f_{x}>n}} \\
		L_{n} \ := \ \set{f:\fa x\in[a+1/n,b],\ex -1/n<h<0,\abs{\Delta_{h}f_{x}>n}} \\
		G_{n} \ := \ \set{f: f \trm{ not monotone on any interval } I\subseteq[a,b] \trm{ of length } 1/n}
	\end{gather*}
	need to check:
	\begin{enumerate}[i)]
		\item $R_{n},L_{n},G_{n}$ open + dense
		\item $\bigcap_{n=1}^{\infty}R_{n}\cap L_{n}\cap G_{n}$ has the required properties
	\end{enumerate}
	second pt is clear; we check the first pt.
	to see $R_{n}$ open, note:
	\begin{equation*}
		\abs{\Delta_{h}f(x)-\Delta_{h}g(x)}\leq\frac{2\norm{f-g}_{\infty}}{h}
	\end{equation*}
	let $f\in R_{n},x\in[a,b-1/n]$. then $\ex h_{x}$ s.t. $\abs{\Delta_{h_{x}}f(x)}>n$.
	since $f$ cts, $\ex T_{x}\ni x$ open intvl and $v_{x}>0$ s.t. $\fa t\in T_{x}$:
	\begin{equation*}
		\abs{\Delta_{h}f(t)}>n+v_{x}
	\end{equation*}
	in particular, by cptness:
	\begin{equation*}
		[a,b-1/n] \ = \ \bigcup_{x\in[a,b-1/n]}T_{x} \ = \ \bigcup_{j=1}^{m}T_{x_{j}}
	\end{equation*}
	let $h=\min(h_{x_{j}})$ and $v=\min(v_{x_{j}})$.
	then, if $\ep<vh/2$ and $\norm{g-f}_{\infty}<\ep$, we have:
	\begin{equation*}
		\abs{\Delta_{h_{x_{j}}}g(t)} \ > \ \abs{\Delta_{h_{x_{j}}}f(t)}-\frac{2\ep}{h_{x_{j}}}
		\ > \ n+v_{x_{j}}-\frac{vh}{h_{x_{j}}} \ > \ n+v-v \ = \ n
	\end{equation*}
	for $t\in T_{x_{j}}$. in particular, $g\in R_{n}$, so $R_{n}$ open.
\end{pf}

